\scanpage{28}
\margin{\figure{\begin{tikzpicture}\draw (0,.2)--(.25,0)--(.55,.35);\end{tikzpicture}}}
\margin{\uline{Soit}\quad Soit \uline{alors}}
\[
\operatorname{Hom}(\overline K,\mathbb C)\xrightarrow{\mathrm{surj}}
\operatorname{Discspec}(\overline K)
\]
est \deleted{\illegible{1}} compatible avec les op. naturelles de $\Gamma$.
\margin{\figure{\begin{tikzpicture}\draw (0,0)--(0,1);\draw (.07,0)--(.07,1);\end{tikzpicture}}}
\deleted{\unsure{Et pour cela}}\inserted{NB} $\Gamma$ est transitif à
\deleted{gauche}\inserted{droite}, et évidemment non transitif à
\deleted{droite}\inserted{gauche}$^*$ (pour que
$\varphi,\varphi':\overline K\rightrightarrows\mathbb C$ soient conjugués par $\Gamma$,
il faut que a) $\varphi|K=\varphi'|K$ et b)
$\varphi(\overline K)=\varphi'(\overline K)$), conditions qui sont des plus
\unsure{draconiennes}\ldots
\margin{$^*$ sauf si $\overline K=\overline{\mathbb Q}$ ($K$ alg. sur $\mathbb Q$
i.e. fini sur $\mathbb Q$)}
\margin{\uline{\unsure{donc les flèches}}\quad\uuline{\unsure{rien}}\quad\uline{\unsure{d'injectif}}}

\uline{\unsure{Autos}} d'une discrétification sur $\overline K$ : on en trouve
d'intéressants par \deleted{\unsure{les}} des
\deleted{\unsure{chemins}}\inserted{``éléments} transcendants'' sur
$K\otimes_{\mathbb Q}\mathbb C$,\ldots\deleted{\unsure{\ldots}}

Est-il possible que $\operatorname{Hom}(\overline K,\mathbb C)\longrightarrow
\operatorname{Discspec}(\overline K)$ soit injectif (i.e bijectif) ?
\deleted{\unsure{Si c'était le cas donc la donnée de}
$\overline K\subset\mathbb C$ \unsure{resp.}}
\inserted{$(K,\overline K,\overline K\xrightarrow{\varphi_{\overline K}}\mathbb C)$}
\deleted{\unsure{resp.}}
i.e. de $(K,\varphi_K:K\mathbin{-}\mathbb C)$, \uline{équivalant} à celle
de $(K,\overline K,\Theta)$ \inserted{$\Theta=\Theta_{\varphi_{\overline K}}$},
\deleted{et tt \unsure{sans autre rigidifications}}
\unsure{Seulement dans le cas trivial} où $\unsure{K|\widehat K}$ est alg. sur $\mathbb Q$.

$\Pi_1\quad\Pi_{1\,\mathrm{ab}}^{\deleted{\unsure{\circ}}}$\quad
\uline{structure de Hodge mixte}
\[
\Pi_{1\,\mathrm{ab}}\otimes_{\mathbb Z}\mathbb C
\]
ou encore structure complexe sur $\unsure{\widetilde\Pi_{1\,\mathrm{ab}}|\Pi_{1\,\mathrm{ab}}}\otimes_{\mathbb Z}\mathbb R$
\mspage{27}
